\documentclass[tikz,border=8pt]{standalone}
\usepackage{tikz}
\usepackage{amsmath}
\usepackage{amssymb}
\usepackage{pgfplots}
\pgfplotsset{compat=1.18}
\usepackage{ctex}
\usetikzlibrary{calc,positioning,arrows.meta}

% LC 169 Majority Element
% Boyer-Moore 投票法：数组 [2,2,1,1,1,2,2]，7步演化

\begin{document}
\begin{tikzpicture}

%% ── 标题 ──────────────────────────────────────────────────
\node[font=\large\bfseries] at (6.0, 0.6)
  {LC 169：Boyer-Moore 投票法（Majority Element）};

%% ── 顶部：原始数组 ────────────────────────────────────────
\foreach \val/\i in {2/0, 2/1, 1/2, 1/3, 1/4, 2/5, 2/6}{
  \node[draw, minimum size=0.75cm, inner sep=0pt,
        fill=white, font=\small\bfseries]
    at (\i*1.0 + 1.5, -0.2) {\val};
  \node[font=\tiny, gray] at (\i*1.0 + 1.5, -0.72) {\i};
}
\node[font=\small,gray] at (0.5,-0.2) {arr:};

%% ── 表格区域 ─────────────────────────────────────────────
% 表头
\def\ty{-1.8}  % 表头 y
\def\rowh{0.65}

\node[font=\small\bfseries,align=center] at (0.7,  \ty) {步};
\node[font=\small\bfseries,align=center] at (2.0,  \ty) {arr[i]};
\node[font=\small\bfseries,align=center] at (3.8,  \ty) {操作};
\node[font=\small\bfseries,align=center] at (6.5,  \ty) {candidate};
\node[font=\small\bfseries,align=center] at (8.8,  \ty) {count};

% 分隔线
\draw[gray!60, thin] (0.1, \ty-0.38) -- (10.0, \ty-0.38);
\draw[gray!60, thin] (0.1, \ty+0.38) -- (10.0, \ty+0.38);

%% 表格数据：7行
% 格式：步号, arr[i], 操作描述, candidate, count, 行高亮色
\def\rows{%
  {0, 2, count=0 \to 初始化 cand=2 count=1, 2, 1, blue!8},%
  {1, 2, cand==arr[i] \to count++,            2, 2, white},%
  {2, 1, cand!=arr[i] \to count--,            2, 1, white},%
  {3, 1, cand!=arr[i] \to count--=0,          2, 0, orange!15},%
  {4, 1, count=0 \to 初始化 cand=1 count=1,   1, 1, blue!8},%
  {5, 2, cand!=arr[i] \to count--,            1, 0, orange!15},%
  {6, 2, count=0 \to 初始化 cand=2 count=1,   2, 1, blue!8}%
}

% 手动写7行（避免 \foreach 嵌套复杂性）

%% 行0
\def\ry{-2.75}
\fill[blue!8] (0.1,\ry-0.32) rectangle (10.0,\ry+0.32);
\node[font=\small] at (0.7, \ry) {0};
\node[font=\small\bfseries, blue!80!black] at (2.0, \ry) {2};
\node[font=\scriptsize, align=center] at (3.8, \ry) {count=0};
\node[font=\scriptsize, align=center] at (5.2, \ry) {初始化};
\node[font=\small\bfseries] at (6.5, \ry) {2};
\node[font=\small\bfseries] at (8.8, \ry) {1};

%% 行1
\def\ry{-3.4}
\node[font=\small] at (0.7, \ry) {1};
\node[font=\small\bfseries, blue!80!black] at (2.0, \ry) {2};
\node[font=\scriptsize, align=center] at (4.3, \ry) {cand==arr[i]，count++};
\node[font=\small\bfseries] at (6.5, \ry) {2};
\node[font=\small\bfseries] at (8.8, \ry) {2};

%% 行2
\def\ry{-4.05}
\node[font=\small] at (0.7, \ry) {2};
\node[font=\small\bfseries, red!70!black] at (2.0, \ry) {1};
\node[font=\scriptsize, align=center] at (4.3, \ry) {cand!=arr[i]，count--};
\node[font=\small\bfseries] at (6.5, \ry) {2};
\node[font=\small\bfseries] at (8.8, \ry) {1};

%% 行3：count归零，高亮
\def\ry{-4.70}
\fill[orange!15] (0.1,\ry-0.32) rectangle (10.0,\ry+0.32);
\node[font=\small] at (0.7, \ry) {3};
\node[font=\small\bfseries, red!70!black] at (2.0, \ry) {1};
\node[font=\scriptsize, align=center] at (4.3, \ry) {cand!=arr[i]，count--=\textbf{0}};
\node[font=\small\bfseries] at (6.5, \ry) {2};
\node[font=\small\bfseries, red!80!black] at (8.8, \ry) {\textbf{0}};

%% 行4：重新初始化，高亮
\def\ry{-5.35}
\fill[blue!8] (0.1,\ry-0.32) rectangle (10.0,\ry+0.32);
\node[font=\small] at (0.7, \ry) {4};
\node[font=\small\bfseries, red!70!black] at (2.0, \ry) {1};
\node[font=\scriptsize, align=center] at (3.8, \ry) {count=0};
\node[font=\scriptsize, align=center] at (5.2, \ry) {初始化};
\node[font=\small\bfseries, red!70!black] at (6.5, \ry) {1};
\node[font=\small\bfseries] at (8.8, \ry) {1};

%% 行5：count归零
\def\ry{-6.00}
\fill[orange!15] (0.1,\ry-0.32) rectangle (10.0,\ry+0.32);
\node[font=\small] at (0.7, \ry) {5};
\node[font=\small\bfseries, blue!80!black] at (2.0, \ry) {2};
\node[font=\scriptsize, align=center] at (4.3, \ry) {cand!=arr[i]，count--=\textbf{0}};
\node[font=\small\bfseries, red!70!black] at (6.5, \ry) {1};
\node[font=\small\bfseries, red!80!black] at (8.8, \ry) {\textbf{0}};

%% 行6：重新初始化，最终结果
\def\ry{-6.65}
\fill[green!12] (0.1,\ry-0.32) rectangle (10.0,\ry+0.32);
\node[font=\small] at (0.7, \ry) {6};
\node[font=\small\bfseries, blue!80!black] at (2.0, \ry) {2};
\node[font=\scriptsize, align=center] at (3.8, \ry) {count=0};
\node[font=\scriptsize, align=center] at (5.2, \ry) {初始化};
\node[font=\small\bfseries, green!50!black] at (6.5, \ry) {\textbf{2}};
\node[font=\small\bfseries] at (8.8, \ry) {1};

% 外框
\draw[gray!50, thin] (0.1, -2.38) rectangle (10.0, -6.97);
% 列分隔线
\foreach \x in {1.2, 2.8, 5.6, 7.6}{
  \draw[gray!30, thin] (\x, -2.38) -- (\x, -6.97);
}
% 行分隔线
\foreach \y in {-3.075,-3.725,-4.375,-5.025,-5.675,-6.325}{
  \draw[gray!20, thin] (0.1,\y) -- (10.0,\y);
}

%% ── 结果标注 ─────────────────────────────────────────────
\node[draw=gray!50,fill=white,thin,rounded corners=2pt,
      font=\small,inner sep=5pt,align=center]
  at (5.0, -7.8) {%
    最终 \textbf{candidate = 2}，即多数元素\\%
    时间 $O(n)$，空间 $O(1)$%
  };

%% ── 算法说明框 ───────────────────────────────────────────
\node[draw=gray!50,fill=white,thin,rounded corners=2pt,
      font=\scriptsize,inner sep=4pt,align=left]
  at (5.0, -9.1) {%
    算法：若 count=0，令 candidate=arr[i]，count=1；\\%
    \hspace{1em}否则 arr[i]==candidate 则 count++，否则 count--
  };

\end{tikzpicture}
\end{document}
